\documentclass[11pt]{article}

\usepackage[T1]{fontenc}
\usepackage[margin=1in]{geometry}
\usepackage{amsmath,amssymb}
\usepackage{booktabs}
\usepackage{hyperref}
\usepackage{xcolor}
\usepackage{enumitem}
\usepackage{listings}

\lstset{
  basicstyle=\ttfamily\small,
  breaklines=true,
  frame=single,
  columns=fullflexible,
  keepspaces=true
}

\title{NR-V2X DeePC Simulation Debugging and Validation Checklist}
\author{}
\date{}

\begin{document}
\maketitle

\section{Purpose}

This document gives a practical validation workflow for the NR-V2X closed-loop DeePC simulation. The goal is to make the simulation easier to debug, easier to reproduce, and easier to defend scientifically.

The central rule is:

\begin{quote}
Do not start debugging with a full 300 s closed-loop run. Validate each layer independently using short tests first.
\end{quote}

\section{Fixed Scenario Definition}

Before running any campaign, freeze one scenario and avoid changing parameters during debugging.

\begin{table}[h]
\centering
\begin{tabular}{ll}
\toprule
Parameter & Recommended value \\
\midrule
Mobility CSV & \texttt{data/processed/ngsim\_us101\_mainline\_active20\_250\_densest\_600s.csv} \\
Simulation time & 300 s for final runs \\
Warmup / cooldown & 10 s / 10 s \\
KPI sample time & 0.5 s \\
KPI window & 1.0 s \\
Control interval & 0.5 s \\
Beacon interval & 0.1 s \\
CAM generation & fixed periodic, ETSI CAM trigger disabled \\
Dynamic active pool & enabled \\
Physical UE pool size & 250 \\
Awareness range & 200 m \\
Controlled input & Tx power only \\
Tx power bounds & 10--23 dBm \\
\bottomrule
\end{tabular}
\end{table}

Do not mix beacon-rate control and Tx-power control until the Tx-power-only pipeline is validated.

\section{Recommended Campaign Folder}

Do not overwrite old results. Use a new campaign folder for each clean restart:

\begin{lstlisting}
data/output/final/ieee_250veh_deepc_campaign_02/
\end{lstlisting}

Recommended structure:

\begin{lstlisting}
commands/
01_training/
02_baselines/
03_closed_loop/
04_analysis/
logs/
\end{lstlisting}

\section{Layered Validation Workflow}

Use the following order:

\begin{enumerate}
    \item validate mobility and active vehicle counts;
    \item run a 10 s fixed-baseline smoke test;
    \item run a 30 s fixed-baseline smoke test;
    \item run a 30 s open-loop training smoke test;
    \item build a small Hankel dataset to test the pipeline;
    \item test the Python DeePC controller using a saved request;
    \item run a 30 s closed-loop smoke test;
    \item plot and validate the closed-loop smoke test;
    \item run full 300 s fixed baseline;
    \item run full 300 s open-loop training;
    \item build final Hankel matrices;
    \item run full 300 s closed-loop DeePC;
    \item compare only matched time windows.
\end{enumerate}

\section{Validation-Only Test}

Before running the radio simulation, run validation-only mode:

\begin{lstlisting}[language=bash]
./ns3 run "nr_v2x_ngsim_deepc_closed_loop_deepc \
  --mobilityCsv=data/processed/ngsim_us101_mainline_active20_250_densest_600s.csv \
  --simTime=300 \
  --warmup=10 \
  --cooldown=10 \
  --sampleTime=0.5 \
  --dynamicActivePool=true \
  --physicalUePoolSize=250 \
  --validateOnly=true"
\end{lstlisting}

Check:

\begin{itemize}
    \item logical vehicle count;
    \item peak simultaneous active vehicles;
    \item peak active core vehicles;
    \item physical UE pool size;
    \item evaluation time window.
\end{itemize}

If validation-only mode fails, do not run the radio simulation.

\section{Required Files Per Run}

Every run should save:

\begin{itemize}
    \item \texttt{command.txt};
    \item \texttt{stdout.log};
    \item \texttt{stderr.log};
    \item \texttt{kpi\_timeseries.csv};
    \item \texttt{kpi\_timeseries\_metadata.json};
    \item \texttt{cbr\_timeseries.csv};
    \item \texttt{plots/};
    \item controller logs for closed-loop runs.
\end{itemize}

The command file is important because many simulation errors come from forgetting which flags were used.

\section{Command Scripts}

Use scripts instead of manual commands. Recommended command scripts:

\begin{lstlisting}
commands/00_validate.sh
commands/01_fixed_10hz_smoke.sh
commands/02_training_smoke.sh
commands/03_build_hankel.sh
commands/04_deepc_smoke.sh
commands/05_fixed_10hz_full.sh
commands/06_training_full.sh
commands/07_deepc_full.sh
commands/08_plot_all.sh
\end{lstlisting}

Manual commands are useful for debugging, but final campaign runs should be script-based.

\section{Dynamic UE Pool Validation}

For the dynamic physical-UE pool, log or verify:

\begin{itemize}
    \item active vehicle count;
    \item active core vehicle count;
    \item assigned vehicle count;
    \item free pool count;
    \item assignment count;
    \item release count;
    \item maximum pool usage.
\end{itemize}

Important invariant:

\[
N_{\mathrm{assigned}} + N_{\mathrm{free}} = N_{\mathrm{pool}}.
\]

Also check:

\[
N_{\mathrm{assigned}} \approx N_{\mathrm{active}},
\]

within the expected active-window logic.

\section{Recommended Pool Assignment Log}

Create or maintain a debug CSV:

\begin{lstlisting}
pool_assignments.csv
\end{lstlisting}

Recommended columns:

\begin{lstlisting}
time_s,event,vehicle_id,node_id,source_l2_id,free_pool_count,active_count
\end{lstlisting}

where event is one of:

\begin{lstlisting}
assign
release
\end{lstlisting}

Check:

\begin{itemize}
    \item no logical vehicle is assigned twice;
    \item no physical node is assigned to two vehicles at once;
    \item a released node returns to the free pool;
    \item source L2 ID is nonzero;
    \item active source L2 IDs do not duplicate.
\end{itemize}

\section{CAM Application State Validation}

When a logical vehicle is released:

\begin{itemize}
    \item CAM dissemination should stop;
    \item the application should be marked unassigned;
    \item the mobility model should park the node outside the road;
    \item the VDP pointer should remain valid and should not dangle.
\end{itemize}

When a new logical vehicle is assigned:

\begin{itemize}
    \item the station ID should match the current logical vehicle ID;
    \item the VDP should point to the current vehicle state;
    \item CAM dissemination should restart;
    \item stale state should not corrupt KPI accounting.
\end{itemize}

Useful debug log:

\begin{lstlisting}
time_s,node_id,event,old_vehicle_id,new_vehicle_id,cam_running
\end{lstlisting}

\section{KPI Sanity Checks}

For every KPI CSV, check:

\begin{align}
0 &\leq \mathrm{PRR} \leq 1, \\
0 &\leq \mathrm{CBR} \leq 1, \\
\mathrm{PIR} &\geq 0, \\
N_{\mathrm{success}} &\leq N_{\mathrm{eligible}}, \\
N_{\mathrm{tx}} &\geq 0.
\end{align}

If any of these fail, the run should not be used.

\section{Plot Early}

After every smoke run, plot:

\begin{itemize}
    \item PRR vs time;
    \item PIR vs time;
    \item CBR vs time;
    \item active vehicle count vs time;
    \item Tx power vs time;
    \item beacon interval vs time;
    \item Tx count vs time;
    \item eligible count vs time;
    \item success count vs time.
\end{itemize}

Do not rely only on mean values. Mean values can hide timing bugs, density effects, and transient errors.

\section{Hankel Dataset Validation}

Before building Hankel matrices, verify:

\begin{itemize}
    \item no NaN values;
    \item no infinite values;
    \item uniform sample time of 0.5 s;
    \item enough rows for the chosen horizons;
    \item input variation is sufficient;
    \item output variation is physically reasonable;
    \item train/validation/test splits are time ordered.
\end{itemize}

For:

\[
T_{\mathrm{ini}} = 20,\qquad N=10,\qquad \Delta t=0.5~\mathrm{s},
\]

the minimum window length is:

\[
T_{\mathrm{ini}} + N = 30~\mathrm{samples}.
\]

In practice, hundreds of samples are needed.

\section{Expected Hankel Dimensions}

For one input:

\[
u = P_{\mathrm{tx}},
\]

three outputs:

\[
y = [\mathrm{PRR},\mathrm{PIR},\mathrm{CBR}]^\top,
\]

and one context signal:

\[
d = N_{\mathrm{active,core}},
\]

the Hankel dimensions should be:

\begin{table}[h]
\centering
\begin{tabular}{ll}
\toprule
Block & Expected dimension \\
\midrule
$U_p$ & $20 \times M$ \\
$U_f$ & $10 \times M$ \\
$Y_p$ & $60 \times M$ \\
$Y_f$ & $30 \times M$ \\
$D_p$ & $20 \times M$ \\
$D_f$ & $10 \times M$ \\
\bottomrule
\end{tabular}
\end{table}

If the context has two columns:

\[
d = [N_{\mathrm{active,core}},T_b]^\top,
\]

then:

\[
D_p = 40 \times M,\qquad D_f = 20 \times M.
\]

\section{Python Controller Validation}

Before connecting to ns-3, test the Python DeePC controller with a saved request.

Check:

\begin{itemize}
    \item OSQP status is \texttt{optimal} or \texttt{optimal\_inaccurate};
    \item Tx power is within $[10,23]$ dBm;
    \item residuals are below $10^{-5}$;
    \item prediction arrays have length 10;
    \item there are no NaN values;
    \item fallback behavior works if OSQP fails.
\end{itemize}

This isolates Python and optimization bugs from ns-3 simulation bugs.

\section{Closed-Loop Controller Timing}

For each closed-loop control step, log:

\begin{itemize}
    \item request time;
    \item response time;
    \item bridge wait time;
    \item solver status;
    \item solve time;
    \item objective value;
    \item equality residuals;
    \item applied Tx power.
\end{itemize}

The controller should solve faster than the control interval:

\[
t_{\mathrm{solve}} < 0.5~\mathrm{s}.
\]

If solve time exceeds the control interval, the closed-loop timing is not realistic.

\section{Fallback Rule}

The closed-loop controller should not crash the simulation when the optimizer fails. Correct fallback behavior is:

\begin{quote}
If OSQP fails, returns non-finite values, or violates bounds, keep the previous Tx power, log the failure, and continue the simulation.
\end{quote}

This makes the simulation robust and easier to debug.

\section{Matched-Window Comparison}

Do not compare runs over different time intervals.

If one run stops at 162 s and another completes 300 s, compare only the common interval:

\[
[t_{\min},t_{\max}] = [10,162]~\mathrm{s}
\]

or another explicitly matched evaluation window.

This is necessary because vehicle density changes over time. Unmatched windows can give misleading mean PRR, PIR, and CBR.

\section{Required Baselines}

Minimum recommended baselines:

\begin{itemize}
    \item fixed 10 Hz, 20 dBm;
    \item fixed 10 Hz, lower Tx power close to the DeePC mean;
    \item fixed 5 Hz, 20 dBm;
    \item open-loop PRBS training run;
    \item closed-loop DeePC.
\end{itemize}

If DeePC reduces mean Tx power, it should be compared against a fixed lower-power baseline. Otherwise, the improvement claim is weak.

\section{Recommended Tools}

Use these tools routinely:

\begin{lstlisting}[language=bash]
rg
git status --short
git diff
./ns3 build <target>
./ns3 run "... --validateOnly=true"
./v2x_env/bin/python -m py_compile scripts/*.py
./v2x_env/bin/python scripts/plot_deepc_bridge_run.py <RUN_DIR>
\end{lstlisting}

For crashes:

\begin{lstlisting}[language=bash]
gdb --args build/scratch/ns3.42-nr_v2x_ngsim_deepc_closed_loop_deepc-default ...
\end{lstlisting}

For memory issues:

\begin{lstlisting}[language=bash]
valgrind --tool=memcheck ...
\end{lstlisting}

For runtime profiling:

\begin{lstlisting}[language=bash]
/usr/bin/time -v ./ns3 run "..."
\end{lstlisting}

\section{Git Practice}

Before major source-code changes:

\begin{lstlisting}[language=bash]
git status --short
git checkout -b campaign-02-clean
\end{lstlisting}

Do not mix unrelated changes in one branch. Keep these separate:

\begin{itemize}
    \item SCI identity modifications;
    \item dynamic UE pool changes;
    \item DeePC formulation changes;
    \item plotting changes;
    \item dataset-building changes.
\end{itemize}

\section{Bottom Line}

The correct simulation strategy is:

\begin{enumerate}
    \item validate each layer independently;
    \item use short smoke tests before full runs;
    \item log every command and metadata file;
    \item check dynamic-pool invariants;
    \item test the controller before connecting it to ns-3;
    \item compare only matched time windows;
    \item change one assumption at a time.
\end{enumerate}

The most important debug artifacts are:

\begin{itemize}
    \item \texttt{pool\_assignments.csv};
    \item \texttt{controller\_diagnostics.csv};
    \item KPI sanity-check script;
    \item matched-window analysis script;
    \item \texttt{command.txt} for every run.
\end{itemize}

Following this workflow prevents long, expensive simulations from being wasted on configuration or accounting errors.

\end{document}
